This appendix supplies the microphysical detail deferred from Sec.~\ref{sec:response}: the tunneling-burst model that turns a per-sensor share into an event train, the single-constant calibration of the count-integral estimator, the Monte-Carlo construction and convergence of the response matrix $R(E_{\rm rec}\,|\,E_{\rm dep})$, and the parameter scan that turns the crossover into a band.

\emph{EMG tunneling-burst model.} On each sensor the trapped-quasiparticle population set by a share $E_s$ drives a two-exponential pulse [Ref.~\cite{ramanathan2026}, Eq.~4], $N_{\rm qp}(t)=\tau_{\rm qp}N_{\rm qp}\,g(t)$ with unit-area shape $g(t)=(e^{-t/\tau_{\rm qp}}-e^{-t/\tau_{\rm inj}})/(\tau_{\rm qp}-\tau_{\rm inj})$, so the instantaneous tunneling rate is $\Gamma_{\rm in}(t)=(K/V_{\rm tr})N_{\rm qp}(t)$. Here $\tau_{\rm qp}$ is the trapped-quasiparticle lifetime, $\tau_{\rm inj}$ the injection time, $K$ the tunneling-rate constant, and $V_{\rm tr}$ the trap volume. The tunneling events form a burst with exponentially-modified-Gaussian (EMG) arrival times, realized with the \texttt{QuasiparticleBurstModel} template of Ref.~\cite{ramanathan2026}; its EMG parameters $(\tau,\mu,\sigma)$ are fixed once per design so that the pdf peak reproduces $\max_t g(t)=p/\tau_{\rm qp}$ (recorded peak factors $p=0.25$ for Al, $0.13$ for Hf). The Poisson mean supplied to the template is the expected number of \emph{tunneling events},
\begin{equation}
n_{\rm ev}=\int_0^\infty \Gamma_{\rm in}(t)\,dt=\frac{K\,\tau_{\rm qp}}{V_{\rm tr}}\,N_{\rm qp},
\label{eq:eventcount}
\end{equation}
a dimensionless count that is \emph{not} the trapped population $N_{\rm qp}$: the two differ by the prefactor $K\tau_{\rm qp}/V_{\rm tr}\simeq0.03$ (Al) and $\simeq0.008$ (Hf). Feeding $N_{\rm qp}$ instead of $n_{\rm ev}$ would mis-scale the saturation onset by these factors, so Eq.~\eqref{eq:eventcount} is enforced everywhere the burst is realized.

\emph{Count-integral calibration.} Each realized event train is censored by the non-paralyzable dead-time model [Eq.~\eqref{eq:censor}] to a per-sensor observed count, and the array sum $N_{\rm obs}$ is mapped to energy by $E_{\rm rec}=C\,N_{\rm obs}$ [Eq.~\eqref{eq:erec}]. The constant $C$ (eV per event) is a \emph{single global} value per design, fixed once on a deeply unsaturated deposit ($E_{\rm dep}=0.01$~eV, no sensor near the ceiling) by requiring $dE_{\rm rec}/dE_{\rm dep}=1$, i.e.\ calibrating the reconstructed energy to estimate the deposited energy; this absorbs the physical yield fraction $\varepsilon$ into $C$, exactly as calibrating against a known line would. Because the unsaturated summed count is exactly proportional to $E_{\rm dep}$, one point fixes $C$; it is never retuned per bin. The recovery of $E_{\rm rec}\approx E_{\rm dep}$ at low energy is therefore a calibration-consistency check (verified to $<1\%$ for both designs), while the deep-saturation plateau is the genuinely predictive content: with censoring disabled the estimator returns exactly $E_{\rm dep}$ at all energies, so the plateau is the modeled bandwidth ceiling and not a numerical artifact.

\emph{Response-matrix Monte Carlo and convergence.} Each $E_{\rm dep}$ column of $R$ is built by histogramming $N_s=5000$ forward realizations of Eq.~\eqref{eq:erec}, on a grid of $584$ deposits sampled uniformly in $\log E_{\rm dep}$ from $10.14$~eV to $197$~MeV ($\simeq$80 bins per decade), the shared grid on which Sec.~\ref{sec:results} folds the deposit spectra. The $\sim\pi r^2$ on-spot sensors are realized individually; the $\sim$10,287 identical off-spot sensors are drawn as one aggregate multinomial ensemble per realization rather than looped. Columns normalize to unity to $10^{-12}$. The per-cell relative Monte-Carlo error is Poisson,
\begin{equation}
\epsilon_{jk}=\frac{1}{\sqrt{N_{jk}}}=\frac{1}{\sqrt{N_s\,R_{jk}}},
\label{eq:mc-converge}
\end{equation}
with $N_{jk}$ the count in cell $(j,k)$; the well-populated (peak) cell reaches $\epsilon\lesssim1.8\%$ at $N_s=5000$ with clean $1/\sqrt{N_s}$ scaling, and every column is reproducible across processes via deterministic \texttt{crc32} sub-seeding. One honesty split is required: where the event count is feasible ($n_{\rm ev}\le2000$) the burst is realized with the full EMG template, but in deep saturation the train would carry $\mathcal{O}(10^8)$ events and is infeasible to materialize, so there the registered count is drawn Poisson about the analytic censored mean. In that regime the relative spread of the summed $E_{\rm rec}$ is $\sim10^{-3}$---the column is effectively a delta---so the approximation is immaterial to $R$; the per-class mean is pinned to the validated analytic estimator in both branches, leaving no kink at the split.

\emph{Crossover-band scan.} The on-spot sensor saturates first, at the per-sensor onset share $E_{\rm onset}$ ($\simeq1.27$~eV for Ta$\to$Al, $\simeq0.77$~eV for Al$\to$Hf); inverting the sharing partition gives the crossover deposit energy $E_{\rm dep}^{\rm cross}$ of Eq.~\eqref{eq:crossover}. Scanning the low-confidence sharing parameters over $f_{\rm prompt}\in[0.1,0.5]$ and $r\in[1,5]$ spreads this from $\sim$8~eV to $\sim$0.93~keV (Ta$\to$Al) and $\sim$4.8~eV to $\sim$0.57~keV (Al$\to$Hf) about the default $\sim$52.9/32.1~eV, the $f_{\rm prompt}$-dominated band quoted in Sec.~\ref{sec:response}. Two limiting anchors bound it: the equal-split limit $f_{\rm prompt}\to0$, where every sensor shares equally and all saturate together, and the whole-array plateau at the default $f_{\rm prompt}$, where even the diffuse share reaches onset,
\begin{align}
E_{\rm dep}^{\rm equal\text{-}split} &= E_{\rm onset}\,N_{\rm sens}, \nonumber\\
E_{\rm dep}^{\rm plateau} &= \frac{E_{\rm onset}\,N_{\rm sens}}{1-f_{\rm prompt}},
\label{eq:onset-band}
\end{align}
with $N_{\rm sens}\simeq10{,}300$. These evaluate to $\sim$13.1~keV (equal-split) and $\sim$18.6~keV (plateau) for Ta$\to$Al, and $\sim$7.9~keV and $\sim$11.3~keV for Al$\to$Hf. The Al$\to$Hf design, with the lower per-sensor onset, saturates before Ta$\to$Al at every point of the scan.
